\documentclass[10pt,a4paper]{amsart}
\usepackage{amsmath,amssymb,amsthm,graphicx}
\usepackage{verbatim}
\newenvironment{reference}{\comment}{\endcomment}
%\newenvironment{reference}{}{}
\newenvironment{slogan}{\comment}{\endcomment}
\newenvironment{history}{\comment}{\endcomment}

% For commutative diagrams we use Xy-pic
\usepackage[all]{xy}

% We use 2cell for 2-commutative diagrams.
\xyoption{2cell}
% Theorem environments.
%
\theoremstyle{plain}
\newtheorem{theorem}[subsection]{Theorem}
\newtheorem{proposition}[subsection]{Proposition}
\newtheorem{lemma}[subsection]{Lemma}

\theoremstyle{definition}
\newtheorem{definition}[subsection]{Definition}
\newtheorem{example}[subsection]{Example}
\newtheorem{exercise}[subsection]{Exercise}
\newtheorem{situation}[subsection]{Situation}

\theoremstyle{remark}
\newtheorem{remark}[subsection]{Remark}
\newtheorem{remarks}[subsection]{Remarks}

\numberwithin{equation}{subsection}

% Macros
%
\def\lim{\mathop{\mathrm{lim}}\nolimits}
\def\colim{\mathop{\mathrm{colim}}\nolimits}
\def\Spec{\mathop{\mathrm{Spec}}}
\def\Hom{\mathop{\mathrm{Hom}}\nolimits}
\def\Ext{\mathop{\mathrm{Ext}}\nolimits}
\def\SheafHom{\mathop{\mathcal{H}\!\mathit{om}}\nolimits}
\def\SheafExt{\mathop{\mathcal{E}\!\mathit{xt}}\nolimits}
\def\Sch{\mathit{Sch}}
\def\Mor{\mathop{\mathrm{Mor}}\nolimits}
\def\Ob{\mathop{\mathrm{Ob}}\nolimits}
\def\Sh{\mathop{\mathit{Sh}}\nolimits}
\def\NL{\mathop{N\!L}\nolimits}
\def\CH{\mathop{\mathrm{CH}}\nolimits}
\def\proetale{{pro\text{-}\acute{e}tale}}
\def\etale{{\acute{e}tale}}
\def\QCoh{\mathit{QCoh}}
\def\Ker{\mathop{\mathrm{Ker}}}
\def\Im{\mathop{\mathrm{Im}}}
\def\Coker{\mathop{\mathrm{Coker}}}
\def\Coim{\mathop{\mathrm{Coim}}}

% Boxtimes
%
\DeclareMathSymbol{\boxtimes}{\mathbin}{AMSa}{"02}

%
% Macros for moduli stacks/spaces
%
\def\QCohstack{\mathcal{QC}\!\mathit{oh}}
\def\Cohstack{\mathcal{C}\!\mathit{oh}}
\def\Spacesstack{\mathcal{S}\!\mathit{paces}}
\def\Quotfunctor{\mathrm{Quot}}
\def\Hilbfunctor{\mathrm{Hilb}}
\def\Curvesstack{\mathcal{C}\!\mathit{urves}}
\def\Polarizedstack{\mathcal{P}\!\mathit{olarized}}
\def\Complexesstack{\mathcal{C}\!\mathit{omplexes}}
% \Pic is the operator that assigns to X its picard group, usage \Pic(X)
% \Picardstack_{X/B} denotes the Picard stack of X over B
% \Picardfunctor_{X/B} denotes the Picard functor of X over B
\def\Pic{\mathop{\mathrm{Pic}}\nolimits}
\def\Picardstack{\mathcal{P}\!\mathit{ic}}
\def\Picardfunctor{\mathrm{Pic}}
\def\Deformationcategory{\mathcal{D}\!\mathit{ef}}
\usepackage{cleveref}
\crefname{originallemma}{original lemma}{original lemmas}
\Crefname{originallemma}{Original lemma}{Original lemmas}
\makeatletter
\newlabel{injectives-lemma-derived-products}{{07D9}{source}{}{sourcepointer-1}{}}
\newlabel{injectives-lemma-derived-products@cref}{{[originallemma][0][]07D9}{[1][0][]source}{}{}{}}
\makeatother
\crefname{originalreference}{original reference}{original references}
\Crefname{originalreference}{Original reference}{Original references}
\makeatletter
\newlabel{lemma-compute-cohomology}{{09A2}{source}{}{sourcepointer-2}{}}
\newlabel{lemma-compute-cohomology@cref}{{[originalreference][0][]09A2}{[1][0][]source}{}{}{}}
\makeatother
\crefname{originalreference}{original reference}{original references}
\Crefname{originalreference}{Original reference}{Original references}
\makeatletter
\newlabel{lemma-w-contractible-hypercovering}{{09A1}{source}{}{sourcepointer-3}{}}
\newlabel{lemma-w-contractible-hypercovering@cref}{{[originalreference][0][]09A1}{[1][0][]source}{}{}{}}
\makeatother
\crefname{originallemma}{original lemma}{original lemmas}
\Crefname{originallemma}{Original lemma}{Original lemmas}
\makeatletter
\newlabel{sites-modules-lemma-sections-over-quasi-compact}{{0935}{source}{}{sourcepointer-4}{}}
\newlabel{sites-modules-lemma-sections-over-quasi-compact@cref}{{[originallemma][0][]0935}{[1][0][]source}{}{}{}}
\makeatother
\title{Cohomology of a quasi-compact and quasi-separated scheme on the pro-etale site commutes with direct sums and homotopy colimits}
\author{The Stacks Project Authors}
\date{stacks-project:proetale.tex:local-snapshot-20261012; GNU FDL 1.2 or later}
\begin{document}
\maketitle
\noindent\textit{Editorial note.} All delivered mathematical text is the author's own, copied from the pinned Stacks Project chapter file proetale.tex; no statement or proof step was rewritten, reordered or removed.\par\medskip
\noindent\textit{Editorial note.} The counted proof is the author's complete argument: the hypercovering computation of both sides as total complexes, the identification of the total complexes through the commutation of sections over quasi-compact objects with direct sums, and the formal derivation for homotopy colimits.\par\medskip
\noindent\textit{Editorial note.} The label lemma-quasi-compact-quasi-separated-commutes-direct-sums names the counted statement, and the proof environment immediately following it in the source is delivered verbatim as its proof.\par\medskip
\noindent\textit{Editorial note.} The counted statement is carried by the author's own lemma environment in the chapter source.\par\medskip
\noindent\textit{Editorial note.} Pre-open gate: python operations/agent-authoring/tools/precard\_gate.py --source proetale.tex --label lemma-quasi-compact-quasi-separated-commutes-direct-sums --start 3329 --end 3369 --local-id 3587 returned RESULT=OK (L3587 free, occupancy FREE, public interval clear). Candidate taken from sources/stacks-proetale/candidates-halfA.json.\par\medskip
\noindent\textit{Editorial note.} The statement and proof are copied from the pinned chapter file; the cited inputs are kept as external pointers with their canonical Stacks tags.\par\medskip
\noindent\textit{Allowed standard prerequisites (editorial).}\begin{itemize}
\item Hypercoverings by affine weakly contractible objects. Same chapter, the lemma producing a hypercovering of X whose every term is a single affine weakly contractible scheme covering X. Supplies the hypercovering K used to compute the cohomology of X as a total complex.
\item Computation of cohomology by a hypercovering. Same chapter, the lemma identifying the cohomology of a bounded below complex on the pro-etale site with the total complex of the alternating double complex of its sections over a hypercovering. Supplies the identification of R-Gamma of each complex with the total complex of the associated double complex.
\item Direct sums of complexes represent direct sums in the derived category. Injectives, the lemma that the termwise direct sum of complexes represents the direct sum in the derived category. Supplies the fact that the termwise direct sum of the chosen representatives represents the direct sum of the objects.
\item Sections over a quasi-compact object commute with direct sums. Modules on Sites, the lemma that taking sections over a quasi-compact object commutes with arbitrary direct sums of sheaves of modules. Supplies the commutation of the sections over the quasi-compact U\_n with the direct sum, which is the heart of the proof.
\end{itemize}\medskip
\ifdefined\hypertarget\hypertarget{sourcepointer-1}{}\fi
\noindent\textit{Source reference.} Original lemma 07D9: sources/stacks-chow/tags.txt, source lines 9404--9404. Prerequisite pointer invoked at proof line(s) 3349; it names a standard fact of a sibling chapter that the counted argument uses, and every step of the counted proof built on it is retained verbatim in the delivered body.\par\medskip
\ifdefined\hypertarget\hypertarget{sourcepointer-2}{}\fi
\noindent\textit{Source reference.} Original reference 09A2: sources/stacks-proetale/pinned-original-source.tex, source lines 3283--3283. Prerequisite pointer invoked at proof line(s) 3350; it names a standard fact of the same chapter that the counted argument uses, and every step of the counted proof built on it is retained verbatim in the delivered body.\par\medskip
\ifdefined\hypertarget\hypertarget{sourcepointer-3}{}\fi
\noindent\textit{Source reference.} Original reference 09A1: sources/stacks-proetale/pinned-original-source.tex, source lines 3216--3216. Prerequisite pointer invoked at proof line(s) 3343; it names a standard fact of the same chapter that the counted argument uses, and every step of the counted proof built on it is retained verbatim in the delivered body.\par\medskip
\ifdefined\hypertarget\hypertarget{sourcepointer-4}{}\fi
\noindent\textit{Source reference.} Original lemma 0935: sources/stacks-chow/tags.txt, source lines 11503--11503. Prerequisite pointer invoked at proof line(s) 3365; it names a standard fact of a sibling chapter that the counted argument uses, and every step of the counted proof built on it is retained verbatim in the delivered body.\par\medskip
\begin{lemma}
\label{lemma-quasi-compact-quasi-separated-commutes-direct-sums}
Let $X$ be a quasi-compact and quasi-separated scheme.
The functor $R\Gamma(X, -) : D^+(X_\proetale) \to D(\textit{Ab})$
commutes with direct sums and homotopy colimits.
\end{lemma}

\begin{proof}
The statement means the following: Suppose we have a family of objects
$E_i$ of $D^+(X_\proetale)$ such that $\bigoplus E_i$ is an object
of $D^+(X_\proetale)$. Then
$R\Gamma(X, \bigoplus E_i) = \bigoplus R\Gamma(X, E_i)$.
To see this choose a hypercovering $K$ of $X$ with $K_n = \{U_n \to X\}$
where $U_n$ is an affine and weakly contractible scheme, see
Lemma \ref{lemma-w-contractible-hypercovering}.
Let $N$ be an integer such that $H^p(E_i) = 0$ for $p < N$.
Choose a complex of abelian sheaves $\mathcal{E}_i^\bullet$
representing $E_i$ with $\mathcal{E}_i^p = 0$ for $p < N$.
The termwise direct sum $\bigoplus \mathcal{E}_i^\bullet$ represents
$\bigoplus E_i$ in $D(X_\proetale)$, see
Injectives, Lemma \ref{injectives-lemma-derived-products}.
By Lemma \ref{lemma-compute-cohomology} we have
$$
R\Gamma(X, \bigoplus E_i) =
\text{Tot}(s((\bigoplus \mathcal{E}^\bullet_i)(K)))
$$
and
$$
R\Gamma(X, E_i) = \text{Tot}(s(\mathcal{E}^\bullet_i(K)))
$$
Since each $U_n$ is quasi-compact we see that
$$
\text{Tot}(s((\bigoplus \mathcal{E}^\bullet_i)(K))) =
\bigoplus \text{Tot}(s(\mathcal{E}^\bullet_i(K)))
$$
by Modules on Sites, Lemma
\ref{sites-modules-lemma-sections-over-quasi-compact}.
The statement on homotopy colimits is a formal consequence of the fact
that $R\Gamma$ is an exact functor of triangulated categories and the
fact (just proved) that it commutes with direct sums.
\end{proof}

\end{document}
